On note $\mathbb{Q}$ l'ensemble des nombres rationnels. Une fonction $f \colon \mathbb{Q} \to \mathbb{Q}$ est
dite \emph{aquesuliane} si elle vérifie la propriété suivante : pour tout $x, y \in \mathbb{Q}$, au moins
une des deux égalités est vérifiée :
\[
f(x + f(y)) = f(x) + y \qquad \text{ou} \qquad f(f(x) + y) = x + f(y).
\]
Démontrer qu'il existe un entier $c$ tel que pour toute fonction aquesuliane $f$, la quantité
$f(r) + f(-r)$ prenne au plus $c$ valeurs distinctes lorsque $r$ parcourt $\mathbb{Q}$, et déterminer la
plus petite valeur possible pour $c$ vérifiant cela.
